% 第 17 章：子类型化的 ML 实现
% Source: source/split/chapters/ch17-ml-subtyping_p212-p215.pdf
% Original print pages: 221--224
% Status: 人工总审中

\chapter{子类型化的 ML 实现}
\label{chap:17}

本章扩展\chapref{10}实现的简单类型 lambda 演算，加入支持子类型化所需的机制，
其中最重要的是检查子类型关系的函数。

\section{语法}
\label{sec:17.1}

类型和项的数据类型定义遵循\cref{fig:15-1,fig:15-3}中的抽象语法。

\begin{taplocamlcode}
type ty =
    TyRecord of (string * ty) list
  | TyTop
  | TyArr of ty * ty

type term =
    TmRecord of info * (string * term) list
  | TmProj of info * term * string
  | TmVar of info * int * int
  | TmAbs of info * string * ty * term
  | TmApp of info * term * term
\end{taplocamlcode}
\taplrustcounterpart{ch17-syntax}
\taplrustsupport{ch17-support}

与纯粹的简单类型 lambda 演算相比，这里新增了类型
\texttt{TyTop}、类型构造子 \texttt{TyRecord}，以及项构造子
\texttt{TmRecord} 和 \texttt{TmProj}。记录及其类型采用最直接的表示：
字段名与对应的项或类型所组成的列表。

\section{子类型关系}
\label{sec:17.2}

\taplsecref{16.1}{sec:16.1}给出的算法子类型关系伪代码可以直接写成下面的
OCaml 函数。

\begin{taplocamlcode}
let rec subtype tyS tyT =
  (=) tyS tyT ||
  match (tyS,tyT) with
    (TyRecord(fS), TyRecord(fT)) ->
      List.for_all
        (fun (li,tyTi) ->
          try let tySi = List.assoc li fS in
              subtype tySi tyTi
          with Not_found -> false)
        fT
  | (_,TyTop) ->
      true
  | (TyArr(tyS1,tyS2),TyArr(tyT1,tyT2)) ->
      (subtype tyT1 tyS1) && (subtype tyS2 tyT2)
  | (_,_) ->
      false
\end{taplocamlcode}
\taplrustcounterpart{ch17-subtype}

这里对伪代码稍作修改，在函数开头增加了反身性检查。运算符 \texttt{(=)}
是普通的相等比较；这里把它写在前缀位置，是因为在另外一些子类型实现中，
它会被替换成对其他比较函数的调用。\texttt{||} 是短路布尔“或”：若第一个
分支得到 \texttt{true}，第二个分支便不会执行。严格说来，这项检查并非必需，
但在真实编译器中却是一项重要优化。现实程序很少实际利用子类型化；换言之，
子类型检查器收到的大多数类型对本来就相等。进一步说，若类型表示能保证结构
同构的类型也具有相同的物理表示——例如在构造类型时使用
\termfirst{哈希合并}{hash consing}（Goto, 1974；Appel 与 Gonçalves,
1993）——这项检查只需一条指令。

记录子类型规则自然需要对列表作一些处理。\texttt{List.for\_all} 把谓词
（第一个参数）应用于列表的每个元素，只有所有应用都返回 \texttt{true} 时
才返回 \texttt{true}。\texttt{List.assoc li fS} 在字段列表 \texttt{fS}
中查找标签 \texttt{li}，并返回相应的字段类型 \texttt{tySi}；若
\texttt{fS} 中不存在该标签，它就抛出 \texttt{Not\_found}，这里捕获该
异常并把结果改为 \texttt{false}。

\section{类型检查}
\label{sec:17.3}

类型检查函数是在先前实现的 \texttt{typeof} 函数上直接扩展而成的。主要
变化出现在应用项的分支：这里要检查实参类型是否为函数所要求参数类型的子类型。
此外，还增加了构造记录和投影字段的两个分支。

\begin{taplocamlcode}
let rec typeof ctx t =
  match t with
    TmRecord(fi, fields) ->
      let fieldtys =
        List.map (fun (li,ti) -> (li, typeof ctx ti)) fields in
      TyRecord(fieldtys)
  | TmProj(fi, t1, l) ->
      (match (typeof ctx t1) with
         TyRecord(fieldtys) ->
           (try List.assoc l fieldtys
            with Not_found -> error fi ("label " ^ l ^ " not found"))
       | _ -> error fi "Expected record type")
  | TmVar(fi,i,_) -> getTypeFromContext fi ctx i
  | TmAbs(fi,x,tyT1,t2) ->
      let ctx' = addbinding ctx x (VarBind(tyT1)) in
      let tyT2 = typeof ctx' t2 in
      TyArr(tyT1, tyT2)
  | TmApp(fi,t1,t2) ->
      let tyT1 = typeof ctx t1 in
      let tyT2 = typeof ctx t2 in
      (match tyT1 with
         TyArr(tyT11,tyT12) ->
           if subtype tyT2 tyT11 then tyT12
           else error fi "parameter type mismatch"
       | _ -> error fi "arrow type expected")
\end{taplocamlcode}
\taplrustcounterpart{ch17-type-of}

记录分支引入了几个此前没有见过的 OCaml 特性。在 \texttt{TmRecord} 分支
中，程序用 \texttt{List.map} 把下面的函数依次应用于每个“字段名／字段项”
二元对，由此把 \texttt{fields} 转换成“字段名／字段类型”列表
\texttt{fieldtys}：

\begin{taplocamlcode}
fun (li,ti) -> (li, typeof ctx ti)
\end{taplocamlcode}
\taplocamlsyntaxonly

在 \texttt{TmProj} 分支中，程序再次使用 \texttt{List.assoc} 查找所选字段
的类型。若查找抛出 \texttt{Not\_found}，程序便报告自己的错误消息。
运算符 \texttt{\^{}} 用于连接字符串。

\begin{exercise}[
  \(\star\star\star\)]
\label{ex:17.3.1}
\taplsecref{16.3}{sec:16.3}说明，在带子类型化的语言中加入布尔值和条件式时，
还需要计算给定类型对之最小公共超类型的辅助函数。命题 16.3.2 的证明给出了
这些算法的数学描述。

\texttt{joinexercise} 类型检查器是不完整的“带子类型化、记录与条件式的
简单类型 lambda 演算”实现。它已经提供基本的解析与打印函数，但
\texttt{typeof} 中缺少 \texttt{TmIf} 分支，也缺少该分支所需的
\texttt{join} 函数。请为该实现加入布尔值、条件式、并与交。
\end{exercise}
\taplauthorsolutionlink{sol:author:17.3.1}{17.3.1}

\begin{exercise}[\(\star\star\)]
\label{ex:17.3.2}
按照\taplsecref{16.4}{sec:16.4}的说明，为 \texttt{rcdsub} 实现加入最小
类型 \(\tapltype{Bot}\)。
\end{exercise}
\taplauthorsolutionlink{sol:author:17.3.2}{17.3.2}

\begin{exercise}[\(\star\star\), \(\nrightarrow\)]
\label{ex:17.3.3}
若应用规则中的子类型检查失败，类型检查器给出的错误消息可能并不能帮助用户
找到问题。即使同时打印预期参数类型与实际实参类型，也可能难以看出差别。例如，
预期类型为
\[
\{x:\{\},y:\{\},z:\{\},a:\{\},b:\{\},c:\{\},d:\{\},
  e:\{\},f:\{\},g:\{\}\}
\]
而实际类型为
\[
\{y:\{\},z:\{\},f:\{\},a:\{\},x:\{\},i:\{\},b:\{\},
  e:\{\},g:\{\},c:\{\},h:\{\}\}
\]
时，很难一眼看出后者缺少字段 \texttt{d}。

可以把 \texttt{subtype} 从“返回 \texttt{true} 或 \texttt{false}”改成
“成功时返回 unit 值 \texttt{()}，失败时抛出异常”。异常会在两个类型真正
无法匹配的位置抛出，因而能够更准确地说明原因。这不会改变类型检查器最终是
接受还是拒绝程序：原来的子类型检查只要返回 \texttt{false}，类型检查器随后
本来就会调用 \texttt{error} 抛出异常。

请重新实现 \texttt{typeof} 和 \texttt{subtype}，使所有错误消息尽可能
提供有用信息。
\end{exercise}
\taplsolutionlink{sol:translator:17.3.3}{17.3.3}

\begin{exercise}[\(\star\star\star\), \(\nrightarrow\)]
\label{ex:17.3.4}
\taplsecref{15.6}{sec:15.6}用从类型推导与子类型推导到纯简单类型 lambda
演算项的翻译，为带记录和子类型化的语言定义了强制转换语义。请实现这些翻译：
修改上面的 \texttt{subtype}，让它构造并返回以项表示的强制转换函数；相应地
修改 \texttt{typeof}，让它同时返回类型与翻译后的项。随后应当求值并打印
翻译后的项，而不是原始输入项。
\end{exercise}
\taplsolutionlink{sol:translator:17.3.4}{17.3.4}
